\documentclass[11pt]{article}

\usepackage[a4paper,margin=1in]{geometry}
\usepackage{amsmath,amssymb}
\usepackage{hyperref}

\setlength{\parindent}{0pt}
\setlength{\parskip}{0.7em}

\title{Bloch-Vector Averaging Across \(N_J\) Sectors}
\author{}
\date{}

\begin{document}
\maketitle

\section*{Purpose}

Use this note whenever changing code that computes, averages, or plots
active-manifold Bloch vectors, Bloch angles, dressed-state directions, or
quantities derived from \(\theta_J,\phi_J\).

\section*{Core Rule}

Do not compare or average raw collective Bloch vectors from different \(N_J\)
sectors as directions.

Raw collective components such as

\begin{equation}
J_x,\qquad J_y,\qquad J_z
\end{equation}

scale with the number of active atoms in the sector. A sector with larger
\(N_J\) can have a larger raw vector simply because it contains more atoms, not
because its single-particle Bloch direction is different. Therefore,
direction-like quantities must be computed from normalized single-particle
components.

For the active manifold \(\{|\downarrow\rangle,|e\rangle\}\), the code uses

\begin{equation}
N_{\rm active}
=
\langle N_\downarrow + N_e\rangle
=
2\bigl(\langle N_e\rangle-\langle J_z\rangle\bigr),
\end{equation}

because

\begin{equation}
J_z=\frac{N_e-N_\downarrow}{2}.
\end{equation}

The normalized active-manifold Bloch direction is then

\begin{equation}
s_x=\frac{2\langle J_x\rangle}{N_{\rm active}},
\qquad
s_y=\frac{2\langle J_y\rangle}{N_{\rm active}},
\qquad
s_z=-\frac{2\langle J_z\rangle}{N_{\rm active}}.
\end{equation}

Angles are computed only after this normalization:

\begin{equation}
\theta=\arccos(s_z),
\qquad
\phi=\operatorname{atan2}(s_y,s_x).
\end{equation}

This is implemented in:

\begin{verbatim}
common.utils.active_manifold_angles(...)
\end{verbatim}

\section*{How Sector Averaging Should Work}

For a direct-sum wavefunction over sectors,

\begin{equation}
|\psi\rangle=\bigoplus_\alpha |\psi_\alpha\rangle,
\end{equation}

where \(\alpha\) may be a homogeneous sector \(N_J\) or an inhomogeneous sector
\((N_{J,1},N_{J,2})\), first compute expectation values with the correct sector
probabilities:

\begin{equation}
\langle J_a\rangle
=
\frac{\sum_\alpha
\langle\psi_\alpha|J_{a,\alpha}|\psi_\alpha\rangle}
{\sum_\alpha \langle\psi_\alpha|\psi_\alpha\rangle},
\end{equation}

and similarly for \(\langle N_e\rangle\). Then compute \(N_{\rm active}\),
\(s_x,s_y,s_z\), \(\theta\), and \(\phi\) from those averaged expectation
values.

Equivalently, this is an active-population-weighted average of the normalized
sector directions. It is not an unweighted average over sectors and it is not an
average of raw collective vector lengths.

Do not do:

\begin{equation}
\theta_{\rm avg}
=
\frac{1}{N_{\rm sectors}}\sum_\alpha \theta_\alpha,
\qquad
\phi_{\rm avg}
=
\frac{1}{N_{\rm sectors}}\sum_\alpha \phi_\alpha.
\end{equation}

Do not do:

\begin{equation}
\mathbf{s}_{\rm avg}
=
\frac{1}{N_{\rm sectors}}\sum_\alpha
\left(
\frac{2J_{x,\alpha}}{N_{{\rm active},\alpha}},
\frac{2J_{y,\alpha}}{N_{{\rm active},\alpha}},
-\frac{2J_{z,\alpha}}{N_{{\rm active},\alpha}}
\right)
\end{equation}

unless an explicitly unweighted diagnostic is desired and clearly labeled as
such.

\section*{Inhomogeneous Group Averages}

For two-group inhomogeneous coupling, group-resolved angles
\((\theta_1,\phi_1)\) and \((\theta_2,\phi_2)\) are computed by applying the same
normalization separately to each group:

\begin{equation}
N_{{\rm active},a}
=
2\bigl(\langle N_{e,a}\rangle-\langle J_{z,a}\rangle\bigr),
\end{equation}

\begin{equation}
s_{x,a}=\frac{2\langle J_{x,a}\rangle}{N_{{\rm active},a}},
\qquad
s_{y,a}=\frac{2\langle J_{y,a}\rangle}{N_{{\rm active},a}},
\qquad
s_{z,a}=-\frac{2\langle J_{z,a}\rangle}{N_{{\rm active},a}}.
\end{equation}

The combined or average angle should be computed by summing the group
expectation values first,

\begin{equation}
\langle J_x\rangle=\langle J_{x,1}\rangle+\langle J_{x,2}\rangle,
\qquad
\langle J_y\rangle=\langle J_{y,1}\rangle+\langle J_{y,2}\rangle,
\end{equation}

\begin{equation}
\langle J_z\rangle=\langle J_{z,1}\rangle+\langle J_{z,2}\rangle,
\qquad
\langle N_e\rangle=\langle N_{e,1}\rangle+\langle N_{e,2}\rangle,
\end{equation}

and then calling \texttt{active\_manifold\_angles(...)} on the summed values.
This avoids arithmetic angle averaging such as \((\theta_1+\theta_2)/2\), which
is generally wrong.

Current implementation reference:

\begin{verbatim}
quantum_trajectories.inhomogeneous_diagnostics.inhomogeneous_group_angles(...)
\end{verbatim}

\section*{Trajectory-Ensemble Averages}

For ensemble observables, average the required expectation values over
trajectories on the common \texttt{t\_eval} grid, then compute the normalized
Bloch direction from the averaged moments.

Current implementation reference:

\begin{verbatim}
quantum_trajectories.aggregator.ensemble_observables(...)
\end{verbatim}

This is the same physical logic used for squeezing: direction and covariance
information should be constructed from ensemble-averaged moments when the target
object is the unconditioned density matrix.

\section*{When To Use This Logic}

Use this logic for:

\begin{itemize}
\item active-manifold angles \(\theta_J,\phi_J\);
\item dressed-state directions such as \(|1\rangle\), \(|c\rangle\),
  \(|j\rangle\), and \(|s\rangle\);
\item plots comparing Bloch angles across sectors, groups, or trajectories;
\item any diagnostic where the intended object is a single-particle direction
  rather than a collective vector length.
\end{itemize}

Do not use this normalization for genuinely extensive observables where the
total collective quantity is the object of interest, such as total
\(\langle J_x\rangle\), total \(\langle N_e\rangle\), total
\(\langle N_J\rangle\), jump rates, or jump counts.

\end{document}
